\documentclass[11pt]{article}
\usepackage[margin=0.75in]{geometry}
\usepackage{hyperref}
\usepackage{amsmath}
\usepackage{setspace}
\usepackage{graphicx}
\usepackage{xcolor}

\hypersetup{colorlinks=true, urlcolor=blue!60!black, citecolor=blue!60!black}
\setstretch{1.1}

\title{\vspace{-1.5cm}\normalsize\textbf{Research Proposal:\\
Rare-Stratum Fidelity in Privacy-Preserving\\
Synthetic Clinical Data}}
\author{
    William Chen\\
    \textit{University of California, Irvine}\\
    Instructor: Prof.\ Sam Behseta\\
    \textit{Summer Session 1}
}
\date{}

\begin{document}
\maketitle
\thispagestyle{empty}

\section*{Introduction}

This project develops synthetic clinical data that preserve the
statistical properties of real data while protecting patient privacy.
These two goals conflict most sharply inside rare strata: uncommon
diagnoses and small comorbidity groups carry little weight in global
fidelity metrics, yet they are often what clinical research cares about
most, and where a generator is most likely to lose structure or
memorize individuals. In an earlier study of prior sensitivity in
Bayesian logistic regression
(\href{https://github.com/ShengPeiWilliam/bayesian-prior-sensitivity}{Bayesian Prior Sensitivity}),
I found that estimate stability was driven by a predictor's
positive-case count, not sample size. This project examines whether
synthetic fidelity behaves the same way at the stratum level.

\section*{Methods}

As a starting point, I would begin from the data and move up to the
tree-based generator, holding the generator fixed and varying only the
input so that any change in fidelity is attributable to the input.

\begin{enumerate}
    \item \textbf{Sweep the positive count}: Hold everything fixed,
    including sample size, and vary only the rare stratum's positive
    count. Compare the structure recovered from the real subset against
    the CART-synthesized data (\texttt{synthpop}) at each count.

    \item \textbf{Sweep the sample size}: Hold the positive count
    fixed and vary the sample size instead, to separate the count
    effect from sample size and prevalence.
\end{enumerate}

I will first validate this on a simulation with a known ground truth,
then extend it to the tree-based generation on MIMIC-IV. All work will
be conducted in R using \texttt{synthpop}, \texttt{rstanarm}, and
\texttt{ggplot2}.

\section*{Goals}

\begin{itemize}
    \item Understand how fidelity and privacy are balanced inside rare
    strata, where preserving structure and protecting individuals pull
    in opposite directions. This tradeoff is the part I most want to
    study.
    \item Compare the two experimental levers directly, sweeping the
    positive count against sweeping the sample size, to see which one
    actually drives rare-stratum fidelity.
    \item Connect this to the tree-based generation used in the
    project, moving from the input data up to the CART-based
    synthesizer and examining where it succeeds or fails on rare
    strata.
\end{itemize}

\newpage

\section*{Preliminary Result}

As a first step, I ran the simulation with a known interaction
(\texttt{chf:lactate}, true value $1.5$) in a rare stratum, sweeping
its positive count with sample size fixed at $500$ (single draw). The
real-data estimate stabilizes near the truth by a count of about $20$;
the synthesized estimate does not recover until roughly $80$.

\begin{figure}[h]
\centering
\includegraphics[width=0.65\textwidth]{../data/figure_saturation.png}
\caption{Recovered \texttt{chf:lactate} interaction versus rare-stratum
positive count, sample size fixed at 500, single draw.}
\end{figure}

\end{document}